\properlane{one-people-through-ages}{One vineyard older and larger than ourselves}
\subsection{The people of God exceed one generation}
Gregory begins with the landowner as Creator and the vineyard as the universal Church. Its extent reaches from Abel to the last elect person at the world's end. The holy persons it brings forth are its branches. The image therefore places the present assembly within a history it did not begin and cannot own. God has continually sent workers to cultivate his people: fathers, teachers of the law, prophets, and apostles. These are not rival employers developing unrelated fields. They serve the one Lord's cultivation.\footnote{Gregory the Great, \textit{Homilies on the Gospels} I.19.1, PL 76, cols. 1154--1155.}

Gregory assigns the hours to successive eras: Adam to Noah, Noah to Abraham, Abraham to Moses, Moses to Christ, and Christ to the world's end. The Gentiles' eleventh-hour arrival belongs to that historical pattern. Yet his vineyard already includes the holy fathers of Israel. The interpretation does not place all earlier people outside salvation and all later people inside it. The ancient faithful are workers in the one vineyard whose reward Christ brings. Its corporate horizon honors the saints of earlier ages while explaining the late gathering of the nations.

Augustine also reads successive historical calls, but he allocates them differently: Abel and Noah belong to the first hour, the patriarchs to the third, Moses and Aaron to the sixth, the prophets to the ninth, and Christians to the eleventh. For both, one Lord gathers his people into promised blessedness across time.\footnote{Augustine, Sermon 87, §5.}

The Sunday readings acquire a corporate depth within this horizon. Isaiah's invitation to return and the psalm's praise are voices of the people already being cultivated. Paul's service and Lydia's reception belong to apostolic gathering. The Entrance addresses a people in need of saving help; the Collect asks for the charity by which its members live together. The Gospel does not authorize the present generation to possess God's goodness as though earlier and later recipients were intruders. It opens the assembly toward a communion larger than its immediate circle.

\subsection{Different ways of understanding the waiting}
Gregory asks the same hard question as Chrysostom: how can those receiving the kingdom murmur? He answers that no murmurer receives heaven and no heavenly saint can murmur. He then interprets the complaint through the delay experienced by the ancient fathers. Though they had lived justly, their entry awaited the Mediator who opened paradise through his death. What the story calls murmuring signifies that long postponement of the promised end.\footnote{Gregory, \textit{Gospels} I.19.4, PL 76, cols. 1156--1157.}

He also illustrates a late arrival rewarded swiftly through the thief who reaches paradise before Peter. The point is not that the thief's entire life was better than Peter's. Nearness to the end can make the interval between calling and entrance very short. A long history of labor does not permit someone to control the moment or fullness of another person's reception. Gregory concludes that even the last place in God's kingdom is cause for joy, which directly undercuts the desire to make salvation valuable chiefly as precedence over another.\footnote{Gregory, \textit{Gospels} I.19.3--4.}

Augustine answers payment order by emphasizing a common resurrection. All receive together, but the first-called have waited much longer. One who receives after a short interval can be described as receiving first relative to the length of waiting, even when the final reception occurs together. Gregory's account stresses Christ opening paradise and swift entrance; Augustine's stresses the shared resurrection and elapsed waiting. Chrysostom, by contrast, treats the complaint as narrative magnification of the latecomer's honor.

For Gregory, Christ brings the earlier faithful and the later called toward the same kingdom. Augustine supplies a parallel account of common blessedness and a further account of how that people is gathered. Neither grants the current assembly the right to treat the history before it as wasted waiting.

\subsection{The workers cultivate one another}
For Gregory, the ecclesial vineyard does not remove individual responsibility. He expressly allows the hours to signify the ages of one person's life. Then he asks whether the hearer has truly begun the Lord's work. A person absorbed only in private advantage has not. Labor consists in charity, devotion, and helping others toward life. Membership is thus a summons to cultivate, not simply a place from which to inspect who else has entered.\footnote{Gregory, \textit{Gospels} I.19.2, PL 76, col. 1155.}

Augustine's Sermon 87 gives this gathering a concrete apostolic movement. The householder's going out signifies Christ becoming known. The risen Lord's mission and Pentecost make the call manifest through preaching. Augustine's final exhortation then turns hearers toward their own friends: concern for those with whom one might live eternally demands active love. His account of conversion includes former persecutors receiving Christ's Blood sacramentally. The people are gathered by a mercy able to transform opposition into membership and service.\footnote{Augustine, Sermon 87, §§9, 14--15.}

That is a substantive addition to Gregory's historical horizon. Gregory explains the vineyard's reach and the work of drawing others toward life; Augustine shows the proclaimed Christ gathering hearers and the sacrament receiving people whose past could have been a ground for exclusion. The call does not merely replace one social group with another. It changes the persons who receive it. Paul's conversion, invoked by Augustine, and Paul's willing service in Philippians belong together: mercy received becomes labor for another's good.

The acclamation's source in Lydia makes the corporate movement visible at household scale. The Lord opens a heart, the message is attended to, and hospitality welcomes the messengers. A universal Church is gathered through particular persons and concrete acts of reception. Paul's willingness to remain supports that same gathering over time. In Aquinas's exposition, conduct worthy of the Gospel extends into concord and cooperation in the clauses following the appointed 27a. Those contextual clauses show the communal direction of the sentence without becoming additional Sunday verses.\footnote{Chrysostom, \textit{Acts} 35; Aquinas, \textit{Super Philippenses} 1, lecture 4, Dessain II, p. 373.}

Isaiah's requirement to abandon an old way challenges inherited membership as surely as it challenges a new convert. The psalm's mercy and nearness give every generation the same reason to return. The Entrance's faithful Lord remains the helper of his people through affliction; the Collect's law of love binds personal devotion to responsibility for a neighbor. These elements keep the historical reading from becoming a diagram in which the listener's only task is to locate an era. The question is whether this member is doing the vineyard's work now.

\subsection{A common table and two Communion voices}
The people's offerings and professed faith lead into a common sacramental participation. Augustine's John exposition gives that communion its costly form. The Shepherd purchased the sheep with his Blood; the preacher did not. Members who receive that gift must be prepared to give themselves for one another. The Church's unity therefore cannot be measured simply by how many people occupy the same space. It is nourished by Christ's gift and expressed through a charity that accepts another member's good as one's own concern.\footnote{Augustine, \textit{John} 47.2.}

If Psalm 119 is used at Communion, Augustine's fuller Latin exposition identifies the speaker as a member or the entire body of Christ. The petition for directed ways can therefore be heard as the body's prayer for faithful action, not merely a private resolution among many unrelated resolutions. Bellarmine's stress on divine assistance reinforces the dependence of each member. If John 10 is used, the good Shepherd's knowledge of his sheep gives the communion its personal center. The wider chapter's gathering into one flock is its context, while the appointed verse names mutual knowledge.

The Prayer after Communion asks that the redemption received become a lived reality in conduct. Within this historical reading, that conduct includes responsibility for those still to be called. A generation that receives the faith holds neither the beginning nor the completion of the vineyard in its own hands. Its opportunity is to be fruitful within the interval given to it: hearing the Word, serving neighbors, and helping others toward the life it hopes to share. The common end makes that labor meaningful; it does not render the workers interchangeable.

\subsection{Four senses of the gathered people}
\begin{fourSenses}
\item[Literal.] One landowner gathers several groups into one vineyard. The appointed sequence includes Israel's prophetic and psalmic witness, Paul's apostolic service, the adapted memory of Lydia's hearing, and Jesus' parable. The Entrance, Collect, offerings, either Communion text, and final prayer address the worship and conduct of a people, with distinct roles in its celebration.

\item[Allegorical.] Gregory identifies the vineyard with the universal Church from Abel to the last elect. Christ brings the fathers and the later called toward one kingdom; Augustine's preaching and Pentecost explain the manifest gathering of the body. The Shepherd purchases and nourishes that body, whose common law is charity.

\item[Moral.] Belonging early gives responsibility for others, not possession of their admission. Gregory's labor of bringing others toward life, Augustine's care for friends, Paul's useful remaining, and Lydia's hospitality make corporate membership active. Obedience sought at Communion becomes a life that cultivates other members rather than resenting their arrival.

\item[Anagogical.] Earlier and later generations are ordered toward the same kingdom. Gregory's account of paradise opened by Christ and Augustine's common resurrection are distinct explanations of the waiting. The Collect's eternal life, Paul's desire for Christ, enduring praise, and the Shepherd's pasture converge on the joy of a gathered people whose last arrivals are welcomed.
\end{fourSenses}
